% SAT Competition benchmark description — cxlb XOR-SLP instances
% Build:  tectonic doc/matmul_cxlb_satcomp.tex
\documentclass[conference]{IEEEtran}
\usepackage{amsmath,amssymb}
\usepackage{booktabs}
\usepackage{url}
\usepackage[hidelinks]{hyperref}
\usepackage{microtype}

\title{Shortest XOR Straight-Line Programs for $3\times 3$ Matrix
  Multiplication:\\ Benchmark Description}
\author{\IEEEauthorblockN{Greg Sidebottom \qquad Claude Fable 5}%
\IEEEauthorblockA{logic repository,
\url{https://github.com/gsidebottom/logic}}}

\begin{document}
\maketitle

\begin{abstract}
We contribute SAT instances deciding ``do $k$ XOR additions
suffice?'' for the output sides of rank-23 bilinear $3\times 3$
matrix-multiplication schemes. Each boundary instance is a live
mathematical question: an UNSAT answer at the right $k$ lower-bounds
the integer additive complexity of a scheme class and, combined with
published exhaustive side analyses, would settle whether 55
additions are possible in known classes (current record: 56, Sun
2026). The family's AND-guarded parity structure defeats current
XOR/Gaussian reasoning while remaining genuinely hard for plain
CDCL: no tested configuration (kissat, CaDiCaL, CryptoMiniSat plain
and native-XOR, Z3 word-level, portfolios with symmetry breaking)
decides the calibrated boundary cells within 10--30 minutes.
\end{abstract}

\section{Problem origin}

A bilinear scheme multiplying two $3\times 3$ matrices with 23
products \cite{laderman} evaluates 23 products of linear forms and
combines them into the 9 output entries. Its \emph{additive
complexity} --- binary $\pm$ additions in a straight-line program,
negation free, no change of basis --- has fallen rapidly: 60
\cite{stapleton}, 59 \cite{mws}, 58 \cite{perminov}, and the current
record \textbf{56} \cite{sun}. Sun's $56 = 13 + 13 + 30$ splits into
two provably optimal \emph{input sides} and an \emph{output side} of
30 additions computing the 9 output forms from the 23 products.

Whether \textbf{55 additions} are possible is open. An exhaustive
orbit-level analysis of all 17{,}376 de Groote classes of the public
Heule--Kauers--Seidl scheme database \cite{hks} (exact input-side
minimization over every representative; research notes in the
repository) reduces the question, for all known classes, to the
output side: \emph{a 55-addition scheme in a known class requires an
output side of $\le 27$ additions on specific representatives} ---
two below the best output count ever observed at this format.

Any $\mathbb{Z}$-coefficient output-side program reduces mod~2 to an
XOR straight-line program, so the decision problem
\begin{quote}
\textbf{SLP($k$):} do $k$ XOR additions suffice to compute the 9
given output forms (vectors in $\mathrm{GF}(2)^{23}$) from the 23
inputs?
\end{quote}
yields sound lower bounds: UNSAT at $k$ proves the integer output
side needs $\ge k{+}1$ additions, and UNSAT at $k=27$ on the right
cells closes entire equivalence classes for 55. These instances
continue the tradition of matrix-multiplication benchmarks in past
competitions \cite{hks}, with a twist of independent solver
interest: the parity constraints are \emph{AND-guarded}, which
prevents CryptoMiniSat's Gaussian elimination from engaging, while
plain CDCL faces genuine parity hardness.

\section{Encoding}

SLP synthesis in the style of Fuhs--Schneider-Kamp \cite{fsk},
simplified by unit-vector inputs. For steps $t = 1,\dots,k$:

\begin{itemize}
\item \textbf{Source selection.} Step $t$ selects exactly two
  sources among the 23 inputs and steps $1..t{-}1$ (sequential
  counter for the exactly-two constraint).
\item \textbf{Values.} Value bits are parity-defined:
  \[
  x_{t,i} \;=\; \mathit{sel}_{t,\mathrm{base}_i} \;\oplus\;
  \bigoplus_{j<t} \bigl(\mathit{sel}_{t,\mathrm{step}_j} \wedge
  x_{j,i}\bigr).
  \]
  Unit-vector inputs make the base contribution a single literal.
  The AND-guarded parities are materialized as Tseitin chains
  (plain CNF) or native \texttt{x}-lines (CryptoMiniSat extension);
  the generator emits both.
\item \textbf{Outputs.} Each of the 9 forms must equal some step
  value (selector-guarded bit equalities; weight-1 forms may match
  an input).
\item \textbf{Symmetry breaking} (optional, default on): all step
  values nonzero; adjacent \emph{independent} steps (step $t{+}1$
  does not consume step $t$) must have strictly lexicographically
  increasing values. Sound: swapping adjacent independent steps
  preserves validity, so every program normalizes by bubble sort,
  and padding above the minimum survives as a dependent chain.
  Dead-step elimination is deliberately \emph{not} used: with it,
  satisfiability is not monotone in $k$ (odd-length padding can
  force a dead step), which corrupts minimum-finding by descent.
\end{itemize}

At $k=29$ with symmetry breaking: 25{,}799 variables, 91{,}426
clauses.

\section{Instances and empirical hardness}

An instance is determined by (output tensor of a scheme
representative, $k$). Tensors come from the public database and the
record chain; the de Groote action turns every class into ${\sim}28$k
distinct cells, so the supply is practically unlimited with $k$ as
the hardness dial:

\begin{itemize}
\item \textbf{SAT phase} ($k \ge$ minimum): minutes at
  minimum${}+1$; witnesses are extracted and replay-verified by the
  generator.
\item \textbf{Deep UNSAT} ($k$ well below minimum): easy.
\item \textbf{Boundary} ($k \in \{\min-1, \min\}$): hard. On the
  calibrated seed cells no tested configuration decides the
  boundary within 10--30 minutes (Apple M4 Pro, one core per
  solver): kissat (600\,s and 1800\,s), CaDiCaL (600\,s),
  CryptoMiniSat~5 (615\,s, plain and native-XOR), Z3 4.16 on a
  word-level QF\_BV formulation (600\,s), and a
  kissat{+}CaDiCaL{+}CMS portfolio with symmetry breaking (900\,s).
\end{itemize}

\begin{table}[h]
\centering
\caption{Calibrated seed cells (SAT witnesses verified).}
{\small
\begin{tabular}{lcc}
\toprule
cell & weight & GF(2) min \\
\midrule
\texttt{sun56} output side (record scheme) & 49 & $\in \{29, 30\}$ \\
\texttt{cn120} output side ($C{=}28$ rep) & 60 & $\in \{27, 28\}$ \\
\bottomrule
\end{tabular}}
\end{table}

Deciding either boundary advances the underlying question directly;
UNSAT at 27 on the fat-sides window cells of the record class would,
combined with the published exhaustions, make ``no 55-addition
scheme exists in the record class'' a theorem with a DRAT
certificate.

\section{Proposed benchmark set}

Twenty instances spanning the phases: the two seed cells at $k \in
\{\min{-}1, \min, \min{+}1\}$; the identity cells of the two other
known 56-addition classes (\texttt{i12w219c23ci},
\texttt{i19w225c4efh}) at their boundaries; and four fat-sides
window cells of the record class at $k=27$; each with and without
symmetry breaking. All are emitted by:

\begin{verbatim}
python3 matmul/cxlb.py --bits SCHEME.bits \
    --k K --dump out.cnf
python3 matmul/cxlb.py --bits SCHEME.bits \
    --k K --dump out.xnf   # native XOR
\end{verbatim}

\noindent(\texttt{--no-sb} disables symmetry breaking; all scheme
bits files are committed in the repository.)

\section{Availability}

Generator, scheme inputs, verification tooling, and research notes:
\url{https://github.com/gsidebottom/logic}. Submitted instances may
be used freely under the competition's standard terms.

\begin{thebibliography}{9}
\bibitem{sun} Y.~Sun. A 56-addition scheme for rank-23 $3\times3$
matrix multiplication. arXiv:2604.27645, 2026.
\bibitem{perminov} D.~Perminov. arXiv:2512.21980, 2025.
\bibitem{mws} E.~M\aa rtensson, P.~Stankovski Wagner, J.~Stapleton.
arXiv:2601.05272, 2025.
\bibitem{stapleton} J.~Stapleton. arXiv:2508.03857, 2025.
\bibitem{hks} M.~J.~H.~Heule, M.~Kauers, J.~Seidl. Local search for
fast matrix multiplication. \emph{SAT 2019}; J.~Symbolic
Computation 104, 2021.
\bibitem{fsk} C.~Fuhs, P.~Schneider-Kamp. Synthesizing shortest
linear straight-line programs over GF(2). \emph{SAT 2010}.
\bibitem{bp} J.~Boyar, R.~Peralta. A new combinational logic
minimization technique with applications to cryptology.
\emph{SEA 2010}.
\bibitem{laderman} J.~Laderman. A noncommutative algorithm for
multiplying $3\times 3$ matrices using 23 multiplications.
Bull.~AMS 82(1), 1976.
\end{thebibliography}

\end{document}
